% 第 29 章　为什么打碎的杯子不会自己复原（完整样章）
\chapter{为什么打碎的杯子不会自己复原}

\step{反常识开场}

请你先在脑子里放一段录像：一只手松开，一只瓷杯下落，砸在地板上，碎成十几片，茶水四溅。现在把这段录像\textbf{倒着放}：四溅的茶水从地板上聚拢回来，碎片从各处飞起，严丝合缝地拼成一只完整的杯子，跳回桌面。

你会觉得倒放版假得可笑。任何观众都能一眼认出“这是倒放”，尽管他可能完全不懂物理。

这个“倒放测试”可以再玩几次。一滴蓝墨水滴进清水，慢慢晕开，整杯水变成均匀的淡蓝色——倒放版是蓝色自动收缩回一滴，一眼假。打开香水瓶盖，几分钟后整个房间都闻得到——倒放版是散布各处的香水分子集体掉头、排队钻回瓶子里，一眼假。滚烫的咖啡放在桌上慢慢变凉——倒放版是凉咖啡从房间里收集热量、自己重新沸腾，一眼假。我们每个人从小就在用这双“一眼假”的眼睛，只是从没问过：我们凭什么认得出来？

现在换一段录像。拍摄一杯水中两个尘埃颗粒的碰撞：一个颗粒慢悠悠飘过来，撞上另一个，两者弹开，各自飘走。把这段录像倒放给任何人看，没有人能分辨哪个是正放、哪个是倒放——两个版本看起来都同样“合理”。再放大胆一点：如果能拍到两个空气分子的对撞，倒放版同样毫无破绽。

这就怪了。杯子由分子组成，尘埃颗粒的碰撞也是分子碰撞的放大版。如果每一个微观的碰撞过程正放倒放都同样合法，为什么亿万个这样的过程加在一起，就出现了鲜明的方向——杯子只会碎，不会拼回去？

你可能已经学过：机械运动服从牛顿定律，而牛顿定律里没有任何一条写着“禁止碎片拼回杯子”。事实上，如果把碎裂瞬间每一个分子的速度精确地反向，牛顿定律会乖乖地把所有碎片送回杯子的形状。力学并不禁止杯子复原。可它从来不发生。一次也没有。

\begin{mainline}
物理学在这里遇到了一条奇怪的裂缝：微观的每一条定律都“不分正反”，宏观的世界却只有一个方向。裂缝中间缺了一样东西——本章要补上的，是\textbf{概率}。
\end{mainline}

\step{先猜一猜}

在继续读之前，请先写下你自己的猜测。下面几种说法，你觉得哪个最接近真相？

第一种猜测：碎片不拼回去，是因为地板把能量“吃掉”了，能量不够，所以复原不了。

第二种猜测：存在某种我们还没学过的力或定律，明文禁止物体从碎变整。

第三种猜测：拼回去并不违法，只是太巧了——需要亿万个分子同时做出恰到好处的运动，而这种巧合的概率小到等于零。

先别急着选。请注意一个事实：碎片拼回杯子，并不违反能量守恒。碎裂时杯子的动能变成了分子热运动和断裂的能量；复原时这些能量再变回杯子的动能和材料的结合能，账面上完全可以平衡。能量守恒定不了方向。那么方向从哪里来？

这个问题不是小题大做。整个十九世纪，最好的物理学家都在它面前摔跤：气体动力学用可逆的分子碰撞解释了一切气体性质，却怎么也解释不了“气体为什么只会扩散不会聚拢”。有人甚至因此拒绝相信原子存在。我们先用双手找一点感觉，再回到这场争论。

另外，请记住你选的是第几种猜测。本章结尾我们要回来对答案——猜错不可耻，可耻的是读完之后忘了自己当初怎么想。物理直觉就是这样一次次被校准的。

\step{动手实验}

\begin{experiment}[摇不回去的豆子]
找一个带盖的透明盒子（保鲜盒就行），再准备两种颜色不同的豆子，比如红豆和绿豆，各二十粒左右。

第一步：把红豆全放在盒子左半边，绿豆全放在右半边，中间不用隔板，只是小心地别碰盒子。盖上盖子，看一眼：两种颜色分得清清楚楚。

第二步：拿起盒子，轻轻摇十下。打开看：豆子混在一起了，红绿相间。

第三步：也是最关键的一步。继续摇，摇一百下、一千下，边摇边观察。豆子有没有可能重新分成“左红右绿”？你可以摇到手酸。历史上没有任何人把摇混的豆子摇回去过。
\end{experiment}

这个实验朴素得近乎可笑，但它和打碎的杯子是\textbf{同一个问题}。豆子没有记忆，不知道什么叫“恢复原状”；盒子里的每次碰撞，正着看倒着看都合法。可混合一旦发生，就再也回不去了。

我们还可以把实验做得更“数字化”。拿四枚硬币，全部正面朝上放在桌上，然后闭上眼睛摇动、撒开，记录正反，重复几十次，数一数：全部正面朝上出现过几次？两正两反出现过几次？如果你想偷懒，也可以用手机里的随机数程序代替：让程序随机出四个 0 或 1，统计“四个都是 0”和“两个 0 两个 1”各出现多少次。结果会稳定地告诉你：两正两反出现的次数远多于全正，大约是它的五六倍。为什么偏偏是这个比例？这就是下一步要拆开的机关。

做实验时顺便回答三个观察题：摇的次数越多，豆子会变得更均匀还是更不均匀？如果只有两粒豆子（一红一绿），“摇回去”还那么难吗？把豆子换成两种细沙，结论变不变？请先写下你的答案，后面的概念会逐个接住它们。

\step{旧模型遇到麻烦}

我们先清点手里的旧模型，看看它们各自在哪里卡壳。

牛顿力学告诉我们：每个豆子的运动由碰撞的力和初速度决定。麻烦在于，力学方程有一种对称性——把时间 \(t\) 换成 \(-t\)，方程照样成立。用第 12 章的语言说：任何一段合法的力学过程，它的“倒放版”也是合法的。一个竖直上抛的小球，倒放就是一个自由下落的小球，两者都严格遵守同一条定律。所以力学无法解释为什么我们从未见过杯子的倒放版。

能量守恒也帮不上忙。前面已经算过账：杯子复原在能量上完全合法。守恒定律只管“总量对不对”，不管“朝哪个方向走”。它像一位只核对总账、从不过问钱花在哪儿的会计。

也许你想说：是摩擦，摩擦生热，热散掉了，所以回不去了。这个说法日常里很好用，但追到底还是不行。热是什么？热就是分子的无规运动。说“摩擦生热所以不可逆”，等于说“分子运动混在一起所以回不来”——这正是我们要解释的现象本身，不是解释。用热解释不可逆，是把答案偷偷塞进了问题里。

还有一条路也容易走进死胡同：怪罪“初始条件太特殊”。有人会说，宇宙碰巧从一个极其特殊的状态出发，所以我们只见单行道。这句话其实说对了一大半——本章结尾的脑洞实验还会回来碰它——但它本身不是解释，只是把问题从“杯子”挪到了“宇宙开局”。而且哪怕不追问宇宙起源，只面对一盒豆子、一杯墨水，方向问题依然原样站着，等一个能算数的回答。

\begin{mainline}
旧模型遇到的麻烦可以压成一句话：\textbf{微观定律允许两个方向，宏观世界却只走一个方向；守恒定律管数量，不管方向。}要走出困境，我们需要的不是一条新力，而是一种新的计数方法。
\end{mainline}

\step{发明新概念}

现在轮到我们自己动手造概念了。回到四枚硬币。

把四枚硬币编上号：1 号、2 号、3 号、4 号。撒一次，你关心的宏观结果可能是“两正两反”。但“两正两反”可以由好几种具体局面实现：1、2 正 3、4 反；1、3 正 2、4 反；1、4 正 2、3 反；2、3 正 1、4 反；2、4 正 1、3 反；3、4 正 1、2 反——一共 6 种。而“全部正面”只有一种实现方式。

\begin{mainline}
我们给这两个层次各起一个名字。你一眼看到、用宏观语言描述的状态（“两正两反”“豆子混匀了”“杯子碎了”）叫\textbf{宏观状态}；每一个成员的确切安排（哪枚硬币朝上、哪个分子在哪个位置以什么速度运动）叫\textbf{微观状态}。一个宏观状态对应的微观状态的个数，记作 \(\Omega\)。
\end{mainline}

四枚硬币的账很好算：“全正” \(\Omega=1\)；“三正一反” \(\Omega=4\)（哪一枚是反面，有 4 种选法）；“两正两反” \(\Omega=6\)；“一正三反” \(\Omega=4\)；“全反” \(\Omega=1\)。加起来 \(1+4+6+4+1=16\)，正好等于 \(2^4\)——四枚硬币一共 16 种微观状态，每种出现的概率相等。这就解释了动手实验里那个“五六倍”：撒出“两正两反”的概率是 \(6/16\)，恰好是“全正” \(1/16\) 的 6 倍。

关键发现出现了：如果每种微观状态的机会均等，那么一个宏观状态出现的概率，就正比于它的 \(\Omega\)。\textbf{不是硬币偏爱混乱，而是“两正两反”这个宏观状态名下挂着更多的微观实现。}你摇豆子时，“左红右绿”的实现数微乎其微，“红绿混杂”的实现数堆积如山——摇匀，只是从一个人烟稀少的宏观状态，随机走进了一个人头攒动的宏观状态。混合不需要任何原因，正如走进大厅不需要任何原因；保持分界才需要奇迹。

把豆子换成气体分子，画面完全一样，只是数字变了。下面这张示意图里，盒子被假想分成左右两半，每个小点是一个分子。左边那幅“全在左半”是低熵状态，右边那幅“均匀分布”是高熵状态——区别不在分子本身（分子一模一样、互不相识），而在宏观安排名下的微观实现数。

\begin{center}
\begin{tikzpicture}[scale=0.85]
  % 左图：分子全在左半
  \draw[thick] (0,0) rectangle (3,2.2);
  \draw[dashed] (1.5,0) -- (1.5,2.2);
  \foreach \px/\py in {0.3/0.4, 0.8/0.9, 0.4/1.5, 1.1/0.5, 0.6/1.9, 1.2/1.3, 0.9/1.7, 0.2/1.1, 1.3/0.9, 0.5/0.8}
    \fill (\px,\py) circle (0.06);
  \node at (1.5,-0.4) {全在左半：实现数极少（低熵）};
  % 右图：均匀分布
  \begin{scope}[xshift=4.6cm]
  \draw[thick] (0,0) rectangle (3,2.2);
  \draw[dashed] (1.5,0) -- (1.5,2.2);
  \foreach \px/\py in {0.3/0.4, 0.8/1.6, 0.4/1.0, 1.1/0.7, 0.6/1.9, 1.9/1.4, 2.4/0.5, 2.7/1.8, 2.1/0.9, 2.6/1.1}
    \fill (\px,\py) circle (0.06);
  \node at (1.5,-0.4) {均匀分布：实现数极多（高熵）};
  \end{scope}
\end{tikzpicture}
\end{center}

顺带交代刚才“大厅”这个比喻的边界。它\textbf{像的部分}：随机乱走的人确实更可能待在大房间里，正如随机碰撞的分子更可能出现在实现数多的宏观状态里。它\textbf{不像的部分}：人有目的、会累、会互相避让，分子没有；房间的大小固定，而宏观状态的“大小”（\(\Omega\)）要靠数出来，不能靠眼睛量。比喻只负责把人领进门，计数才负责把话说准。

十九世纪的玻尔兹曼把这个想法推到底。他提出：给每个宏观状态贴一个数，用来度量它名下微观状态的多少，这个数叫\textbf{熵}，记作 \(S\)。微观状态越多，熵越大。一个孤立系统自发演化的方向，就是走向名下微观状态更多的宏观状态——不是因为有什么力在推，而是因为随机乱走的系统，几乎必然走进那个最大的房间。这个想法在当年激烈到让玻尔兹曼备受攻击：许多同行拒绝接受“看不见的原子”和“统计出来的方向”。今天，原子的照片和涨落的直接测量早已判定了这场争论的胜负。

这里必须立刻说清一件常被讲错的事。流行的说法是“熵就是混乱程度”。这个比喻\textbf{像的部分}是：整齐的局面（全正、左红右绿）确实实现数少，杂乱的局面的确实现数多。\textbf{不像的部分}有三处。第一，“混乱”是人的主观感受，而 \(\Omega\) 是可以老老实实数出来的客观数字，同一个房间对整洁的妈妈和散漫的孩子“乱度”不同，实现数却只有一个。第二，有些看起来极其整齐的结构恰恰是高熵的帮手——冬天的雪花从水汽里结晶出来，冰晶整齐得令人惊叹，但结晶过程把热释放给周围环境，水分子加环境整体的实现数反而增加了。第三，“乱不乱”取决于你用什么宏观语言描述：一副牌对新出厂的定义是乱的，对赌场洗牌标准可能还不够乱。用“看起来乱不乱”判断熵，常常会被眼睛骗到。本章只用实现数说话。

\step{一句话模型}

\begin{oneline}
打碎的杯子不会复原，不是因为定律禁止，而是因为“复原”名下的微观状态少得可怜——系统随机演化，几乎必然走向实现方式更多的宏观状态；熵就是实现方式多少的度量。
\end{oneline}

这句话值得拆开读三遍。第一遍读“不是定律禁止”：方向不是禁令，是可逆定律和计数共同长出来的结果。第二遍读“随机演化”：系统没有目标，它只是乱走，是大数把乱走的平均效果压成了确定的方向。第三遍读“度量”：熵不是标签，是有单位、能测量、能参与计算的量，下一章它会直接走进热机的效率公式。

\step{数学翻译器}

现在把“多”和“少”翻译成数字，你会发现事情比想象的极端得多。

先看硬币怎么数。\(N\) 枚硬币，“恰好 \(k\) 枚正面”的微观状态数，是从 \(N\) 枚里挑 \(k\) 枚当正面的选法数，数学上记作
\[
\Omega(N,k)=\binom{N}{k}=\frac{N!}{k!\,(N-k)!}.
\]
四枚硬币：\(\Omega(4,0)=1\)，\(\Omega(4,2)=6\)，\(\Omega(4,4)=1\)，和前面数的完全一致。总微观状态数是
\[
\sum_{k=0}^{N}\binom{N}{k}=2^{N}.
\]

检查极端情形：\(N=4\) 时 \(2^4=16\)，对上了。\(k=N\)（全正）时 \(\Omega=1\)，也对——全正只有一种摆法。再看中间：\(N\) 是偶数时，\(\Omega\) 在 \(k=N/2\) 处取最大值，也就是“恰好各半”是实现数最多的宏观状态。

下面这张图把计数结果画出来：横轴是宏观状态（四枚硬币中正面的枚数），纵轴是该宏观状态的实现数 \(\Omega\)。中间鼓、两头塌。

\begin{center}
\begin{tikzpicture}[scale=0.9]
  \draw[-{Stealth}] (0,0) -- (6.4,0) node[right] {正面枚数 $k$};
  \draw[-{Stealth}] (0,0) -- (0,3.4) node[above] {$\Omega$};
  \fill[blue!40] (0.55,0) rectangle (1.05,0.25);
  \fill[blue!40] (1.95,0) rectangle (2.45,1.0);
  \fill[blue!40] (3.35,0) rectangle (3.85,1.5);
  \fill[blue!40] (4.75,0) rectangle (5.25,1.0);
  \fill[blue!40] (5.75,0) rectangle (6.25,0.25);
  \node at (0.8,0.5) {\small 1};
  \node at (2.2,1.25) {\small 4};
  \node at (3.6,1.75) {\small 6};
  \node at (5.0,1.25) {\small 4};
  \node at (6.0,0.5) {\small 1};
  \foreach \px/\kk in {0.8/0, 2.2/1, 3.6/2, 5.0/3, 6.0/4}
    \node at (\px,-0.35) {\small $\kk$};
\end{tikzpicture}
\end{center}

请特别注意这个形状随 \(N\) 的变化。\(N=4\) 时，中间只比两头高几倍；\(N=100\) 时，“恰好 50 正”的实现数约为 \(10^{29}\)，而全正仍然只有 1 种；正面比例偏离各半越远，实现数按指数暴跌。硬币越多，分布图越瘦，最后瘦成一根立在正中央的针。\textbf{微观状态均等这条假设，经过组合计数，被放大成了宏观上几乎铁一般的规律。}

现在放大到一百枚硬币。全部正面的概率是 \(1/2^{100}\)。\(2^{100}\) 大约是 \(1.27\times10^{30}\)。假如你每秒撒一次硬币，从宇宙大爆炸（约 \(1.4\times10^{10}\) 年前，合 \(4.4\times10^{17}\) 秒）撒到今天，大约撒了 \(4\times10^{17}\) 次，见到全正的机会仍然只有千亿分之一量级。而“接近各半”（比如正面在 45 到 55 之间）的概率超过七成。撒一百枚硬币，你其实是在围观一次小型的时间之箭。

体会一下 \(2^N\) 的增长有多凶：十枚硬币，全正的概率约千分之一，碰运气能见到；二十枚，约百万分之一，一生难得一见；四十枚，约万亿分之一，全人类一起撒也指望不上。每多一枚硬币，极端状态的机会就砍掉一半，而分子世界里 \(N\) 不是几十，是 \(10^{23}\) 量级。不可逆不是什么深刻的禁令，只是指数增长碾过去之后留下的平地。

\begin{mathbridge}
真正的飞跃发生在取对数这一步。粒子数一大，\(\Omega\) 动辄是 \(10^{10^{23}}\) 这样的怪物，直接写没法写。玻尔兹曼的主意是取自然对数，把天文数字压成平常的数：
\[
S = k_{\mathrm{B}} \ln \Omega,
\]
其中 \(k_{\mathrm{B}}\approx1.38\times10^{-23}\ \mathrm{J/K}\)，叫玻尔兹曼常数，负责把“个数”换算成能和温度、能量对接的单位（焦耳每开尔文）。

这个式子刻在玻尔兹曼的墓碑上。它为什么这样写？三个理由。第一，\(\ln\) 把乘法变加法：两个独立系统放在一起，总实现数 \(\Omega=\Omega_1\Omega_2\)，取对数后 \(S=S_1+S_2\)，熵像能量一样可以相加。第二，\(\ln\Omega\) 随 \(\Omega\) 单调增长，保住了“实现数越多熵越大”的顺序。第三，用 \(k_{\mathrm{B}}\) 做系数后，这个从“数数”得来的熵，和热学里从热量与温度定义的熵完全一致——这在第 30 章会看到。

先给单位一点直觉。焦耳每开尔文，意思是“每升高一度要搭进去多少能量级别的账”。把一勺热水倒进冷水，整个过程的熵变只有 \(10^{-3}\) 焦耳每开尔文量级——看起来很小，但别忘了它是拿 \(10^{-23}\) 这么小的常数乘出来的，反推回去，\(\Omega\) 翻的倍数是 \(e\) 的 \(10^{20}\) 次方量级。日常一勺水的混合，背后是微观实现数翻了无法书写的倍数。

什么时候用它？系统近似孤立、微观状态可以近似看作机会均等时。什么时候不能乱用？系统很小（几个粒子）、或者粒子间有强关联时，涨落会大到让“几乎必然”失效，后面边界一节专门讲。
\end{mathbridge}

再看一个真实尺度的估算。房间里的一摩尔空气，约 \(6\times10^{23}\) 个分子。问：某个瞬间，这些分子碰巧全部挤在房间左半边、右半边空无一物的概率有多大？每个分子在左在右各半机会，全部在左的实现数占总实现数的比例是
\[
\frac{1}{2^{6\times10^{23}}}.
\]
这个数小到什么程度？把它写成十进制，小数点后的零的个数大约是 \(1.8\times10^{23}\) 个——把这些零排成一行，每个零占一毫米，能从地球排到太阳再排回来上亿次。与之相比，“杯子碎片自己拼回去”要求的巧合，量级与此相当甚至更甚：不仅位置要对，每个分子的速度方向也要对。日常经验里“从不发生”，换成数字就是“不可能到这种程度”。

\begin{challenge}
“分子几乎均匀分布在两边”不只是一个概率偏大的说法，它有精确的数学形状。设左半房间有 \(n\) 个分子，\(n\) 服从二项分布，平均值是 \(N/2\)，涨落（标准差）是 \(\sqrt{N}/2\)。注意相对涨落 \(\sqrt{N}/N=1/\sqrt{N}\)：\(N\) 越大，相对涨落越小。一摩尔分子，相对涨落约 \(1.3\times10^{-12}\)——均匀到十一位小数。这就是为什么宏观世界看起来“确定”，尽管底层全是概率：不是概率消失了，而是大数把涨落压没了。

用斯特林近似 \(\ln N!\approx N\ln N-N\)，还可以算出 \(N\) 枚硬币“恰好各半”的实现数 \(\binom{N}{N/2}\approx 2^{N}\sqrt{2/(\pi N)}\)。取对数：\(\ln\Omega\approx N\ln 2\)。于是这堆硬币的熵 \(S\approx N k_{\mathrm{B}}\ln 2\)——熵正比于粒子数。这个“熵是广延量”的性质，正是热力学计算一切平衡态的起点。
\end{challenge}

\step{模型的边界}

每一个模型都要交代三件事：假设了什么、在哪里成立、在哪里会被误用。本章的模型也不例外。

\textbf{假设一：微观状态机会均等。}整个推理的地基是“每种微观安排概率相同”。这个假设在宏观孤立系统里被无数次验证（它还有一个名字，叫等概率原理），但它终究是假设，不是从牛顿定律推出来的定理。什么保证它成立、系统要多“乱撞”多久才算达到均等，是统计物理深处的难题，第 31 章还会继续追。

\textbf{假设二：系统要足够大。}“几乎必然走向大 \(\Omega\) 状态”这句话里的“几乎”，威力来自天文数字般的粒子数。把系统缩小，铁律就软化成概率。三枚硬币撒一次，全正的概率是 \(1/8\)——并不罕见。所以“摇豆子摇不回去”成立，“摇三枚硬币摇不回全正”不成立。这不是理论出错，而是理论诚实：它预言了小系统里不可逆会破功。

一个漂亮的反直觉反例来自花粉。1827 年植物学家布朗看到水中花粉颗粒无规则地乱动，后来明白那是水分子从各方向撞出来的。如果水中悬浮的颗粒足够小、时间足够短，分子撞击的涨落可以真的“推它一把”做一点正功——在微米和纳秒尺度上，“热变功”天天发生。现代实验室里，物理学家已经能用微米级的小珠直接测出这些涨落，结果和统计预言严丝合缝。热力学第二定律从来不说“绝对禁止”，它说的是“宏观上几乎必然”。把统计规律当成和能量守恒同级的绝对禁令，是对它最常见的误用。

\textbf{假设三：系统是孤立的。}开放系统完全可以局部“逆流”——只要向环境付出足够的代价。冰箱把水冻成整齐的冰，局部实现数大幅下降；代价是冰箱背面的散热管把更多的热排进厨房，整间屋子（水加空气）的总实现数增加了。生命是更壮观的例子：一个受精卵长成结构精密的人，局部的熵降得惊人，但人一辈子要吃、要呼吸、要散热，身体和环境合在一起的总熵只增不减。\textbf{“生命违反熵增”是伪命题，因为它偷偷换了系统的边界。}

最后单列一类滥用。一旦“熵”成了流行词，它就被拖去解释一切：“房间不收拾所以熵增”“信息爆炸所以社会熵增”“人心散了所以团队熵增”。这些说法的问题不在结论对错，而在于它们引用的根本不是那个可以数 \(\Omega\) 的物理量。物理学里的熵有单位、有数值、能放进方程参与计算；修辞里的“熵”只是“乱”的洋气的说法。分清这两者的办法很简单：问一句“你说的熵，怎么测？”

不过，“熵是数出来的”这个核心想法确实走出了热学，而且走得很远。香农建立信息论时，要度量“一条消息包含多少不确定性”，最后写下的公式和玻尔兹曼熵在数学上是同一个形状——因为信息量本质上也是“可能情形的个数”的对数，这是第 32 章的主角。更惊人的是黑洞：物理学家发现黑洞的熵正比于它表面积，一个黑洞名下挂着的微观实现数，比同样体积里任何其他东西的都要多——宇宙的“最大房间”竟然是黑洞。从四枚硬币到黑洞，同一把计数的钥匙开了跨度四十个数量级的门。

\begin{mainline}
本章模型的使用说明：孤立、宏观、可以计数时，熵增方向几乎必然；系统小，涨落可以翻盘；系统开放，局部可以减熵，但环境付账更多；离开这三个条件谈熵，多半是在打比方，不是在做物理。
\end{mainline}

\step{回头解释开场现象}

现在回到那两段录像。

尘埃颗粒的碰撞，涉及少量分子的一次事件。把每个参与者的速度反演，得到的倒放过程同样是一次合法碰撞，实现它的微观状态并不比正放少。所以你分不出正反——这个尺度的物理真的没有方向。

杯子的碎裂则是另一回事。碎裂后的世界（碎片的位置、分子的热运动方式）名下挂着多到无法想象的微观状态；而“碎片飞回、拼合、跳上桌面”这个宏观状态，要求亿亿个分子在同一瞬间各自处在精确到苛刻的位置、朝着精确到苛刻的方向运动，名下挂着的微观状态相对少得可怜。概率之比大约又是那种“零的个数能排到太阳”的数字。\textbf{倒放版不是违法，是离谱}——它离谱到宇宙这辈子都等不到一次。

墨水、香水、咖啡同理。晕开的墨水名下微观状态极多，聚拢回一滴名下的极少；所以前者天天发生，后者从未发生。你那双“一眼认出倒放”的眼睛，其实是亿万年演化训练出来的概率探测器。

这也解释了一个细节：摔杯子时声音四散、碎片发烫。机械能变成了分子的无规热运动，而热运动正是实现数最多的能量存放方式——能量一旦“撒”进热里，就像豆子撒进了盒子，再也聚拢不回来。这和第 21 章讲的能量耗散是同一件事的两副面孔：那边从能量流看，这边从实现数看。能量没有消失，只是住进了最大的那个房间。

动手实验留下的三个观察题，现在也能回答了。摇得越久，豆子越均匀，因为均匀状态名下的实现数最大，随机乱走只会越陷越深，不存在“摇过头又分开”。只有两粒豆子时，“摇回去”轻而易举——每种安排的机会相当，全在左边的概率高达 \(1/4\)，这正是“系统要足够大”那条边界在提醒你。换成细沙结论不变，因为不可逆不挑材料，只数实现数；沙粒之间的摩擦会改变混合的快慢，却改变不了方向。

至于开头的直觉猜测，现在可以对答案了。猜测一（能量不够）错在：复原的能量账是平的，方向不由能量守恒决定。猜测二（存在禁令）错在：没有任何定律明文禁止复原，牛顿定律甚至欢迎倒放。猜测三对：方向来自概率的悬殊——这就是时间之箭在宏观世界浮现的方式。微观定律没有箭头，箭头是大数定律用天文数字的悬殊“统计”出来的。

\step{脑洞实验与挑战题}

\begin{thought}[等一次全正]
假如宇宙是一台永不停止的抛硬币机器，那么只要等得够久，再离谱的巧合也会发生。数学上有庞加莱回归定理：一个孤立力学系统，只要等足够长的时间，会任意接近地回到它当初的微观状态——打碎的杯子、散尽的热量、混合的豆子，原则上都会“重演”。

但请看这个时间尺度。对一个只有房间大小、约 \(10^{26}\) 个粒子的系统，回归时间的数量级大约是 \(10^{10^{25}}\) 年（注意是指数塔）。宇宙年龄只有 \(1.4\times10^{10}\) 年。把宇宙年龄当作一秒，回归时间远远超过“把全宇宙每个原子都变成一个宇宙、每个那样的宇宙再活一次”所需的全部时间。所以回归定理和日常经验不冲突：它说“原则上会回来”，经验说“你永远等不到”——两句话都对，隔着的是指数。
\end{thought}

\begin{thought}[妖精要把守的门]
想象盒子里装着混合均匀的两种气体分子，中间隔板上有一扇无摩擦的小门，门边守着一个小妖精，见到快分子从左边来就开门放它去右边，见到慢分子从右边来也放它去左边。久而久之，右边越来越热，左边越来越冷——没用任何能量，熵似乎下降了。这个“麦克斯韦妖”困扰了物理学家一百多年。问题出在哪？提示：妖精不是白干活的，它要测量、要记忆、要擦除记忆。后来人们发现，擦除一比特记忆本身就要向环境放热、产生熵，账一算平——妖精的“信息”原来也是物理的。这个故事把本章和第 32 章的信息、第 30 章的热力学第二定律缝在了一起。
\end{thought}

\begin{puzzles}
\item \textbf{洗牌的账。}一副新扑克牌按花色点数排得整整齐齐，你彻底洗一次牌，它就变成看起来杂乱无章的顺序。有人说：“洗牌让熵增加了。”请用本章的语言判断这句话对不对，并说明：为什么“整齐”只有一种（或很少几种）排列，而“杂乱”有天文数字种？如果洗一次牌偶然洗回了出厂顺序，这违反任何物理定律吗？

  \textit{思路提示：} 52 张牌的排列总数是 \(52!\approx 8\times10^{67}\)。“整齐”这个宏观状态名下只有极少的微观排列，“杂乱”名下挂着几乎全部排列。洗回出厂顺序概率约为 \(10^{-68}\)，不违反任何定律——只是你永远抽不到这个签。注意：牌局算不算熵，取决于你约定的宏观描述方式，这正是“混乱程度”说法靠不住的地方。

\item \textbf{压瘪的乒乓球。}把压瘪的乒乓球丢进热水，它会自己鼓起来。球内的气体分子做了什么？为什么热水里的分子不会“商量好”一起躲开凹陷处、让球保持瘪的？

  \textit{思路提示：} 热水把能量传给球内气体，分子运动更剧烈，各方向撞击球壁的机会均等；“躲开凹陷”要求大量分子的速度方向出现特定关联，那样的微观状态在总数中占比小到可以忽略。均匀撞击是实现数最大的安排。

\item \textbf{两个房间，一个结局。}A 房间先打开香水喷一下，B 房间把同样一瓶香水喷在室外大风天里。哪个房间的气味更快“散开”？两种情况里，香水分子的微观定律有区别吗？那是什么决定了散开的快慢？

  \textit{思路提示：} 微观定律完全一样。快慢的区别在于宏观条件：室外有定向气流和更大的空间，分子可去的实现数增长得更快。熵增的方向相同，速率却可以差几个量级——本章只保证方向，不保证时刻表。

\item \textbf{反向的冰箱。}有人说：“冰箱证明熵可以减少，所以热力学第二定律是错的。”请分别指出这句话里被偷换的两个概念，并各用一句话纠正。

  \textit{思路提示：} 被偷换的是“系统的边界”（冰箱内部不是孤立系统，必须连同散热的环境一起算）和“熵增的性质”（第二定律是统计性的几乎必然，不是绝对禁令，但冰箱根本不构成挑战——它每秒都在用更大的环境熵增买内部的熵减）。
\end{puzzles}

本章从一个假得可笑的倒放视频出发，找到了藏在所有宏观不可逆现象背后的同一把钥匙：不是力，不是禁令，而是计数——微观实现的个数之差，悬殊到超出任何日常尺度，于是概率在宏观世界里凝固成了方向。我们也顺带还清了一笔旧账：第 21 章里“能量耗散”只说能量变得不好用，却没说为什么“不好用”是条单行道；现在知道了，单行道就是熵增。

下一章我们要把这把钥匙插进热力学的锁孔：既然方向由熵增给出，那么“热为什么只从高温流向低温”“热机效率为什么有上限”“绝对零度为什么到不了”，都将是同一个答案的不同说法。再往后，第 31 章会把“数微观状态”发展成一整套统计物理，第 32 章则追问信息和熵到底有什么关系。碎杯子这道日常小事，门口排着的队伍还很长。
